We find the same broad ordering of components with the mixed model. Its descriptive comparison assigns approximately 61\% to the shared field pattern, 29\% to university-specific field differences and 11\% to the residual, compared with 65\%, 28\% and 7\% from nested least squares. The rounded mixed-model percentages need not sum to 100\%. The estimated variances are $\sigma_a^2=0.027$, $\sigma_b^2=0.039$ and $\sigma^2=0.015$ squared grade points, and the REML fit converged with L-BFGS. The interaction variance exceeds the residual variance, supporting a university profile beyond unexplained annual fluctuations. This does not establish persistence by itself; we test that interpretation on later years in Section~\ref{sec:res-prediction}. The two percentage splits are not identical estimands: one uses reductions in residual sums of squares after university-year adjustment, while the other combines fitted fixed-field variation with estimated interaction and residual variances. We saved the partially pooled profiles and conditional uncertainty intervals in \texttt{results/mixed\_model\_profiles.csv}, so the comparison need not rely on raw university averages alone.
